% Ad Familiares 11.7 (ad-familiares-11-07) --- English translation
% Translated by Claude (Anthropic) from the Latin (Perseus).
% Style guide: STYLE.md.

\ciceroLetterOpener{M. CICERO S. D. D. BRVTO IMP. COS. DESIG.}

\ciceroSection{1} When Lupus brought me, Libo, and your cousin Servius together at my house, what my own view was you have, I believe, learned from M. Seius, who was present at our conversation; the rest, although Graeceius followed close on Seius's heels, you will be able to learn from Graeceius all the same.

\ciceroSection{2} But the main point, which I want you to take in most carefully and to keep in mind, is this: in the preservation of the liberty and welfare of the Roman people, do not wait for the authority of a Senate that is not yet free --- lest you both condemn your own act (for it was by no public decision that you freed the commonwealth, which makes that achievement all the greater and more illustrious), and judge the young man, or rather the boy, Caesar to have acted rashly in undertaking so great a public cause on private judgement, and finally judge that men of the countryside, but bravest of men and the best of citizens, were out of their minds --- first the veteran soldiers, your old comrades-in-arms, then the Martian Legion, the Fourth Legion, which judged their own consul a public enemy and threw themselves into defending the welfare of the commonwealth. The will of the Senate ought to be taken for its authority when its authority is hindered by fear.

\ciceroSection{3} Lastly, the cause you have undertaken is now of such a kind that nothing can be left half-done --- first by the Ides of March, and lately by the raising of a new army and the assembling of forces. For which reason you ought to be so prepared, so resolved, for everything, that you do nothing merely because you are ordered, but rather carry through deeds that will be praised by all with the highest admiration.
